% =====================================================================
\section{Профессиональная диагностика: Guru Meditation}
% =====================================================================
\lessonintro
  {превратить пугающий лог падения в конкретную строку и предотвратить повторение ошибки}
  {все предыдущие уроки, особенно задачи, память и сеть}
  {у устройства есть понятный журнал состояния и путь к поиску редких сбоев}
  {Serial Monitor; файл собранной прошивки ELF для декодирования адресов}

\subsection{Паника --- это полезный отчёт}

Когда ESP-IDF обнаруживает критическую ошибку, он печатает \code{Guru Meditation Error}, причину,
регистры и backtrace, а затем обычно перезапускает чип. Не удаляйте этот лог: скопируйте его целиком
вместе со строкой версии прошивки. Самые частые причины для начинающего проекта --- обращение по
нулевому указателю, выход за границы массива, переполнение стека задачи и использование объекта после
\code{free()}.

\begin{center}\small
\begin{tabularx}{\textwidth}{@{}l X X@{}}
\toprule
Симптом & Вероятная причина & Первое действие \\
\midrule
\code{LoadProhibited} & нулевой/испорченный указатель & проверить указатель и место создания объекта \\
\code{Stack canary} & стек задачи мал & измерить High Water Mark, увеличить стек \\
WDT timeout & задача/ISR не отдаёт CPU & убрать долгую работу, добавить ожидание \\
Brownout detector & питание просело & проверить источник, конденсаторы и нагрузку \\
\bottomrule
\end{tabularx}
\end{center}

\subsection{Как читать backtrace}

Адреса в backtrace сами по себе не говорят номер строки. После каждой сборки IDE создаёт файл
\code{.elf}, в котором есть соответствие адресов исходному коду. ESP Exception Decoder или утилита
\code{xtensa-esp32s3-elf-addr2line} переводят адреса в имена функций и номера строк \emph{только для
той же сборки}. Не декодируйте лог от старой прошивки новым ELF: результат будет правдоподобным, но
ложным.

Рабочий порядок такой:
\begin{enumerate}
  \item сохраните полный лог, версию, причину сброса и шаги, после которых он случился;
  \item найдите первую пользовательскую функцию в декодированном стеке, а не последнюю системную;
  \item добавьте проверку входных данных и краткие диагностические сообщения вокруг места ошибки;
  \item повторите сценарий много раз и только затем считайте исправление подтверждённым.
\end{enumerate}

\subsection{Разбор настоящего сбоя}

Программа ниже обращается к объекту, который так и не был создан:

\codecap{Скетч с ошибкой: указатель \code{sensor} равен \code{nullptr}}
\codeinput{l14_panic}

Через секунду после запуска в Serial Monitor появится отчёт (сокращён):

\begin{Verbatim}[fontsize=\scriptsize,bgcolor=codebg]
Guru Meditation Error: Core  1 panic'ed (LoadProhibited). Exception was unhandled.
Core  1 register dump:
PC      : 0x42002a1c  PS      : 0x00060830  A0      : 0x82002b45  A1      : 0x3fcebd80
...
EXCVADDR: 0x00000000  LBEG    : 0x40056f5c  LEND    : 0x40056f72  LCOUNT  : 0x00000000

Backtrace: 0x42002a19:0x3fcebd80 0x42002b42:0x3fcebda0 0x4200593e:0x3fcebdc0
Rebooting...
\end{Verbatim}

\begin{itemize}
  \item \code{LoadProhibited} --- чтение из запрещённого адреса; \code{EXCVADDR: 0x00000000} уточняет:
        читали по нулевому адресу, то есть через нулевой указатель.
  \item \code{Core 1} --- сбой случился на ядре 1, где работает \code{loop()}.
  \item \code{Backtrace} после декодирования ELF-файлом той же сборки превращается в строки кода:
\end{itemize}

\begin{Verbatim}[fontsize=\scriptsize,bgcolor=codebg]
0x42002a19: readSensor(Sensor*) at sketch.ino:5
0x42002b42: loop() at sketch.ino:11
0x4200593e: loopTask(void*) at cores/esp32/main.cpp:74
\end{Verbatim}

Первая пользовательская функция в стеке --- \code{readSensor()}, строка 5. Исправление: создать объект
(\code{Sensor s\{1, 0\}; sensor = \&s;}) и проверять указатель перед использованием.


\subsection{Watchdog: защита от «всё зависло»}

Watchdog перезагружает ядро, если задача не даёт работать idle-задаче. В \code{loop()} и задачах
должны быть точки, где вы отдаёте процессор: \code{vTaskDelay()}, ожидание очереди с тайм-аутом или
короткая неблокирующая работа. В ISR допустимы лишь установка флага, счётчик или отправка события
через API \code{...FromISR}. \code{delay()}, Serial, сеть и выделение памяти внутри ISR --- типичный
путь к WDT.

\subsection{Питание --- тоже часть отладки}

Если лог сообщает Brownout, код может быть совершенно исправен. Wi-Fi при передаче создаёт короткие
пики тока; слабый USB-кабель, маленькая батарея или питание сервопривода от 3V3 платы дают провал.
Сначала отключите мощную нагрузку, измерьте напряжение возле ESP32 и объедините земли источников.
Мотор, реле и сервопривод питаются от отдельного подходящего источника через ключ/драйвер, а не GPIO.

\subsection{Минимальный журнал здоровья}

Раз в минуту полезно печатать: время работы, причину последнего сброса, свободную RAM,
\code{WiFi.RSSI()} и запас стека важных задач. Не печатайте эти значения из ISR и не превращайте
диагностику в бесконечный поток: один компактный отчёт помогает увидеть медленную утечку памяти и
падение качества сети без вреда для основной функции устройства.

\begin{summary}
\begin{itemize}[leftmargin=*]
  \item Полный лог + ELF той же сборки превращают адреса backtrace в строку исходного кода.
  \item WDT означает блокировку CPU, а Brownout часто означает проблему питания, не алгоритма.
  \item Диагностируйте гипотезу воспроизводимым сценарием и проверяйте исправление повторными запусками.
\end{itemize}
\end{summary}

\begin{selfcheck}
\begin{enumerate}
  \item Почему для декодирования backtrace нужен ELF-файл именно той же сборки?
  \item О чём говорит Brownout и где искать его причину?
  \item Что нужно делать в задачах, чтобы watchdog не перезагружал плату?
\end{enumerate}
\end{selfcheck}

\begin{tasks}
\begin{enumerate}
\item \lv{1} Выведите причину последнего сброса и расшифруйте её своими словами.
\item \lv{1} Составьте шаблон отчёта об ошибке: версия, причина сброса, лог, шаги, питание и периферия.
\item \lv{2} Намеренно вызовите ошибку (деление на нуль указателя) и разберите backtrace по строкам.
\item \lv{2} Включите предупреждения компилятора и найдите хотя бы одно предупреждение в своём скетче.
\item \lv{3} Добавьте в скетч урока 5 раз в минуту строку здоровья: время, куча, RSSI и причина сброса.
\item \lv{4} Добавьте сторожевой таймер и проверьте, что намеренно зависшая задача приводит к перезагрузке, а причина сохраняется в журнал.
\end{enumerate}
\end{tasks}
